\chapter{Részletes tervek}
\comment{A dokumentum célja, hogy pontosan specifikálja az implementálandó osztályokat, beleértve a privát attribútumokat és metódusokat, ezek definícióját is.A dokumentum második fele részletesen be kell mutassa a korábban definiált be-és kimeneti nyelv szintakszisát felhasználva, hogy mely tesztekkel lesz a prototípus ellenőrizve.}
\section{Módosítások az előző dokumentumhoz képest}
\subsection{Bemeneti nyelv}
\begin{itemize}
%INICIALIZÁLÓS PARANCSOK
    \item nehezseg [ertek]
    \begin{itemize}
        \item Leírás: Beállítja a játék nehézségét.
        \item \texttt{[érték]}: easy, medium, hard
    \end{itemize}

    \item letrehoz [tipus] [nev]
    \begin{itemize}
        \item Leírás: Létrehoz egy adott objektumot: típusa lehet TakarítóJátékos, BuszosJátékos, Hókotró (maximum egyet lehet létrehozni kezdésként), Busz, Autó
        \item \texttt{[tipus]}: TakarítóJátékos, BuszosJátékos, Hókotró, Busz, Autó
        \item \texttt{[nev]}: egyedi azonosító, amely a későbbiekben a parancsokban használható
    \end{itemize}

    \item szerkeszt [nev] [ujnev] 
    \begin{itemize}
        \item Leírás: Változtatja a korábban létrehozott objektum nevét.
        \item \texttt{[nev]}: a szerkesztendő objektum neve
        \item \texttt{[ujnev]}: az új név, amely egyedi azonosító kell, hogy legyen
    \end{itemize}

    \item kivalaszt [nev]
    \begin{itemize}
        \item Leírás: Kiválasztja a megadott nevű objektumot, amelyre a további parancsok hatni fognak.
        \item \texttt{[nev]}: A kiválasztandó objektum neve
    \end{itemize}

    \item betolt [fajlnev]
    \begin{itemize}
        \item Leírás: Betölti a megadott fájlból az objektumokat.
        \item \texttt{[fajlnev]}: A betöltendő fájl neve
    \end{itemize}

    %PÁLYA ÉS TÉRKÉP PARANCSOK
    \item utletrehoz [nev] [cs1] [cs2] [hossz] [tipus] [savszam]
    \begin{itemize}
        \item Leírás: Létrehoz egy új utat a pályán.
        \item \texttt{[nev]}: Az út neve.
        \item \texttt{[cs1]}: Az utat kezdő csomópontja.
        \item \texttt{[cs2]}: Az utat végző csomópontja.
        \item \texttt{[hossz]}: Az út hossza.
        \item \texttt{[tipus]}: Az út típusa: alagut, normal, hid
        \item \texttt{[savszam]}: Az út sávainak száma.
    \end{itemize}

    \item lerakodas [utnev] [sav] [tipus] [ertek]
    \begin{itemize}
        \item Leírás: Megadja, hogy mi van lerakódva egy adott út adott sávján. 
        \item \texttt{[utnev]}: Az út neve, amelyen a lerakódás történik.
        \item \texttt{[sav]}: Az út sávjának azonosítója.
        \item \texttt{[tipus]}: Ho, Jeg, So, Feltortjeg, Zuzalek
        \item \texttt{[ertek]}: A lerakódott tárgy értéke.
    \end{itemize}

    %MOZGATÓ PARANCSOK
    \item lepes [ut] [sav]
    \begin{itemize}
        \item Leírás: Egy jármű mozgatása a pályán, egy szomszédos útra. 
        \item \texttt{[ut]} : Az út neve amire szeretne lépni a játékos, amelynek a jelenlegi út szomszédságában kell lennie.
        \item \texttt{[sav]} : Opcionális paraméter, amely megadja, hogy a jármű melyik sávba szeretne lépni az adott úton, amennyiben az út több sávos.
    \end{itemize}

    \item wait [nev]
    \begin{itemize}
        \item Leírás: Lépés kihagyása, a játékos nem hajt végre lépést.
    \end{itemize}

    %Menedzsment parancsok

    \item hokotrovasarlas [nev]
    \begin{itemize}
        \item Leírás: A játékos megvásárol egy új hókotró járművet, amennyiben telephelyen tartózkodik és rendelkezik a megfelelő összeggel.
        \item \texttt{[nev]}: a megvásárolni kívánt hókotró neve, amely egyedi azonosító kell, hogy legyen
    \end{itemize}

    \item allapot [nev]
    \begin{itemize}
        \item Leírás: Kiírja egy adott objektum (játékos, hókotró, út) aktuális tulajdonságait.
        \item \texttt{[nev]}: az állapot lekérdezni kívánt objektum neve. "Mind" esetén mindegyiket kiírja. 
    \end{itemize}
    
    \item terkep
    \begin{itemize}
        \item Leírás: Szövegesen kilistázza a csomópontokat és a köztük lévő utakat, hogy lehessen tudni, honnan hova lehet lépni.
    \end{itemize}
    
%HÓKOTRÓS PARANCSOK
    \item fejcsere [hokotro] [ujfej]
    \begin{itemize}
        \item Leírás: Kicseréli a hókotró fejét egy új fejre, amennyiben a játékos rendelkezik a megfelelő fejjel a készletében, és telephelyen tartózkodik.
        \item \texttt{[hokotro]}: a jelenlegi hókotró, amin fejet szeretnénk cserélni
        \item \texttt{[ujfej]}: a kívánt új fej típusa, amire cserélni szeretnénk a jelenlegi fejet
    \end{itemize}

    \item fejvasarlas [hokotro] [fejtipus]
    \begin{itemize}
        \item Leírás: Új fejet vásárol a játékos a készletére, amennyiben telephelyen tartózkodik és rendelkezik a megfelelő összeggel.
        \item \texttt{[hokotro]}: a jelenlegi hókotró, amelyhez fejvásárlást szeretnénk végrehajtani
        \item \texttt{[fejtipus]}: a kívánt fej típusa, amelyet vásárolni szeretnénk
    \end{itemize}

    \item sotoltes [hokotro]
    \begin{itemize}
        \item Leírás: Amennyiben a játékos telephelyen tartózkodik, feltölti a hókotró sókészletét maxra. 
        \item \texttt{[hokotro]}: a jelenlegi hókotró, amelynek sókészletét szeretnénk feltölteni
    \end{itemize}

    \item kerozintoltes [hokotro]
    \begin{itemize}
        \item Leírás: Amennyiben a játékos telephelyen tartózkodik, feltölti a hókotró kerozinkészletét maxra.
        \item \texttt{[hokotro]}: a jelenlegi hókotró, amelynek kerozinkészletét szeretnénk feltölteni
    \end{itemize}

    \item zuzalektoltes [hokotro]
    \begin{itemize}
        \item Leírás: Amennyiben a játékos telephelyen tartózkodik, feltölti a hókotró zuzalékkészletét maxra.
        \item \texttt{[hokotro]}: a jelenlegi hókotró, amelynek zuzalékkészletét szeretnénk feltölteni
    \end{itemize}

    \item takarit [hokotro] [sav]
    \begin{itemize}
        \item Leírás: A hókotro vévigmegy két csomópont között, és takarít az adott sávon. 
        \item \texttt{[hokotro]}: a jelenlegi hókotró, amelyik takarítani fog
        \item \texttt{[sav]}: az út azon sávja, amit takarítani tervezünk
    \end{itemize}

%RENDSZER PARANCSOK
    \item havazas 
    \begin{itemize}
        \item Leírás: Havazás eseményét indítja el, amelynek hatására a pályán hó jelenik meg. 
    \end{itemize}

    \item kilepes
    \begin{itemize}
        \item Leírás: Kilép a játékból
    \end{itemize}

    \item help
    \begin{itemize}
        \item Leírás: Megjeleníti a parancsok listáját
    \end{itemize}


\end{itemize}
Példák a bemeneti nyelvre:
\begin{verbatim}
    utletrehoz a cs1 cs2 50 alagut 1
    utletrehoz b cs2 cs3 50 normal 1
    lerakodas b 1 ho 10
    utletrehoz c cs3 cs4 50 normal 1   lerakodas c 1 feltortjeg 10
    utletrehoz d cs4 cs5 50 normal 1
    lerakodas d 1 jeg 5
    utletrehoz e cs4 cs5 50 normal 1
    lerakodas e 1 zuzalek 5
    lerakodas e 1 jeg 5
    utletrehoz f cs4 cs5 50 normal 1
    lerakodas e 1 zuzalek 5
    lerakodas e 1 ho 10

    Letrehoz Hokotro H1
    Kivalaszt H1
    Lepes a
    Lepes b
    Lepes c
    Lepes e
    Allapot mind
\end{verbatim}

\section{Osztályok és metódusok tervei}


\subsection{Jarmu}
\begin{class-template-responsibility}
    Absztrakt ősosztály az összes mozgó entitás számára. Kezeli a baleset miatti ideiglenes mozgásképtelenséget, amely az utat is eltorlaszolja
\end{class-template-responsibility}

\begin{class-template-attribute}
    \classitem{\#name: String}{A jármű neve/azonosítója.}
    \classitem{\#owner: String}{A jármű tulajdonos játékosának neve.}
    \classitem{\#currentUt: String}{Az aktuális út neve; \texttt{null} esetén a jármű telephelyen van.}
    \classitem{\#savIndex: int}{Az aktuális sáv indexe az adott úton.}
    \classitem{\#disabledTime: int}{A hátralévő mozgásképtelen körök száma. 0 vagy kisebb érték esetén a jármű mozgásképes.}
\end{class-template-attribute}
\begin{class-template-method}
    \classitem{-Jarmu(name: String)}{Privát konstruktor. Inicializálja az alapállapotot: \texttt{name := bemenet}, \texttt{currentUt := null}, \texttt{savIndex := 0}, \texttt{disabledTime := 0}.}
    \classitem{+name(): String \textit{(NamedEntity)}}{Visszaadja a jármű azonosító nevét.}
    \classitem{+renameTo(newName: String): void \textit{(NamedEntity)}}{Frissíti a jármű nevét a megadott új névre.}
    \classitem{+type(): String \textit{(NamedEntity)}}{Alapértelmezett típusnév: \texttt{"Jarmu"}. Leszármazottak felüldefiniálhatják.}
    \classitem{+statusLine(state: GameState): String \textit{(NamedEntity)}}{A jármű állapotának szöveges reprezentációját adja vissza (pozíció + állapot).}
    \classitem{+canMove(): boolean}{Mozgásképesség ellenőrzése. Igaz, ha \texttt{disabledTime <= 0}, egyébként hamis.}
    \classitem{+tickTime(): void}{Körléptetés. Ha \texttt{disabledTime > 0}, eggyel csökkenti az értéket.}
    \classitem{+vegrehajtLepes(target: Ut, targetSav: int, state: GameState): void}{Lépést hajt végre a célútra/célsávba. Ellenőrzi a szomszédosságot a jelenlegi úthoz képest, frissíti a pozíciót, majd lefuttatja az esetleges baleseti logikát és a célút-elérési hookot.}
    \classitem{\#onCelUtElerve(target: Ut, state: GameState): void}{Védett kiterjesztési pont. Alapértelmezésben nem végez műveletet; a leszármazottak (pl. \texttt{Busz}) felüldefiniálhatják célút-eléréskor.}
    \classitem{-maybeCrash(sav: Sav, state: GameState): void}{Privát segédmetódus. Nem hókotró járműnél a sáv jégviszonyai és nehézség alapján balesetet sorsol; siker esetén \texttt{disabledTime := 2}, balesetszámlálót növel és eseményt naplóz.}
\end{class-template-method}



\subsection{Auto}
\begin{class-template-responsibility}
    Önműködő (NPC) jármű, amely az otthona és a munkahelye között próbál eljutni a legrövidebb úton. Lépéseivel, súlyát felhasználva hozzájárul a hó jéggé tömörítéséhez. Kiemelt felelőssége a várakozási logika: ha az út elzáródik, nem tervez újra, nem fordul meg, hanem a legutolsó pontnál feltorlódva várakozik, amíg a hókotrók fel nem szabadítják az utat.\end{class-template-responsibility}
\begin{class-template-baseclass}
    Jarmu \baseclass Auto
\end{class-template-baseclass}
\begin{class-template-attribute}
    \classitem{(nincs saját attribútum)}{Az Auto osztály nem deklarál új adattagot, minden állapotát a \texttt{Jarmu} osztályból örökli.}
    \classitem{\#name: String (örökölt)}{A jármű azonosító neve.}
    \classitem{\#owner: String (örökölt)}{A tulajdonos játékos neve.}
    \classitem{\#currentUt: String (örökölt)}{Az aktuális út neve; \texttt{null} esetén telephelyen van.}
    \classitem{\#savIndex: int (örökölt)}{Az aktuális sáv indexe az úton.}
    \classitem{\#disabledTime: int (örökölt)}{A még hátralévő mozgásképtelen körök száma.}
\end{class-template-attribute}
\begin{class-template-method}
    \classitem{-Auto(name: String)}{Privát konstruktor. Inicializálja a \texttt{Jarmu} ősosztály állapotát: \texttt{name = bemenet}, \texttt{currentUt = null}, \texttt{savIndex = 0}, \texttt{disabledTime = 0}. Az \texttt{Auto} példány létrehozása a parancskezelőn keresztül történik.}
    \classitem{+name(): String \textit{(örökölt)}}{Visszaadja az azonosító nevet.}
    \classitem{+renameTo(newName: String): void \textit{(örökölt)}}{Frissíti az objektum nevét. A hívó oldali logika felel azért, hogy az új név egyedi legyen.}
    \classitem{+statusLine(state: GameState): String \textit{(örökölt)}}{Szöveges állapotképet ad vissza. A metódus determinisztikusan képezi le a belső állapotot emberileg olvasható formára (pozíció + állapot + típus + név).}
    \classitem{+canMove(): boolean \textit{(örökölt)}}{Mozgásképesség vizsgálata. Igazat ad vissza, ha \texttt{disabledTime <= 0}, különben hamisat.}
    \classitem{+tickTime(): void \textit{(örökölt)}}{Idő frissítés. Ha \texttt{disabledTime > 0}, akkor eggyel csökkenti; egyébként változatlanul hagyja.}
\end{class-template-method}

\subsection{Busz}
\begin{class-template-responsibility}
    Megadott állomások között ingázó, a buszsofőr menedzser által irányított jármű. Feladata a pénzteremtés a megtett körök alapján. Az autókhoz hasonlóan súlyával tömöríti a havat. Baleset esetén mozgásképtelenné válik és eltorlaszolja az utat. \end{class-template-responsibility}
\begin{class-template-baseclass}
    Object \baseclass Jarmu \baseclass Busz
\end{class-template-baseclass}
\begin{class-template-attribute}
    \classitem{-completedTrips: int}{A busz által teljesített fordulók száma. Kezdeti értéke 0, és a rendszer növeli, amikor a busz végállomáshoz kapcsolódó útra lép.}
    \classitem{\#name: String (örökölt)}{A jármű azonosító neve.}
    \classitem{\#owner: String (örökölt)}{A tulajdonos játékos neve.}
    \classitem{\#currentUt: String (örökölt)}{Az aktuális út neve; \texttt{null} esetén telephelyen van.}
    \classitem{\#savIndex: int (örökölt)}{Az aktuális sáv indexe az úton.}
    \classitem{\#disabledTime: int (örökölt)}{A még hátralévő mozgásképtelen körök száma.}
\end{class-template-attribute}
\begin{class-template-method}
    \classitem{-Busz(name: String)}{Privát konstruktor. Meghívja a \texttt{Jarmu} konstruktorát a név inicializálásához, majd \texttt{completedTrips := 0}.}
    \classitem{+name(): String \textit{(örökölt)}}{Visszaadja az azonosító nevet.}
    \classitem{+renameTo(newName: String): void \textit{(örökölt)}}{Frissíti az objektum nevét. Az egyediség ellenőrzése a hívó logika feladata.}
    \classitem{+statusLine(state: GameState): String \textit{(örökölt)}}{A busz szöveges állapotképét adja vissza (pozíció és mozgásképesség).}
    \classitem{+canMove(): boolean \textit{(örökölt)}}{Igazat ad vissza, ha a busz nem mozgásképtelen (\texttt{disabledTime <= 0}).}
    \classitem{+tickTime(): void \textit{(örökölt)}}{Idő frissítés. Ha \texttt{disabledTime > 0}, akkor eggyel csökkenti; egyébként változatlanul hagyja.}
    \classitem{\#onCelUtElerve(target: Ut, state: GameState): void \textit{(Override)}}{Célút-eléréskor, ha az út végállomáshoz kapcsolódik, növeli a \texttt{completedTrips} értékét, majd a tulajdonos játékosnak +40 pénzt jóváír és eseményt naplóz.}
\end{class-template-method}



\subsection{Hokotro}
\begin{class-template-responsibility}
    A takarító menedzser által mozgatott jármű, amely az aktuálisan rászerelt fej segítségével tisztítja az utakat. A hókotró különleges tulajdonsága, hogy fizikailag immunis az akadályokra. 
\end{class-template-responsibility}
\begin{class-template-baseclass}
    Object \baseclass Jarmu \baseclass Hokotro
\end{class-template-baseclass}
\begin{class-template-attribute}
    \classitem{-aktivFej: Fej}{Az aktuálisan felszerelt takarítófej. Kezdetben \texttt{SoproFej}.}
    \classitem{-so: int}{A sókészlet mennyisége.}
    \classitem{-kerozin: int}{A kerozinkészlet mennyisége.}
    \classitem{-zuzottko: int}{A zúzottkő készlet mennyisége.}
    \classitem{\#name: String (örökölt)}{A jármű azonosító neve.}
    \classitem{\#owner: String (örökölt)}{A tulajdonos játékos neve.}
    \classitem{\#currentUt: String (örökölt)}{Az aktuális út neve; \texttt{null} esetén telephelyen van.}
    \classitem{\#savIndex: int (örökölt)}{Az aktuális sáv indexe az úton.}
    \classitem{\#disabledTime: int (örökölt)}{A még hátralévő mozgásképtelen körök száma.}
\end{class-template-attribute}
\begin{class-template-method}
    \classitem{-Hokotro(name: String)}{Privát konstruktor. Meghívja a \texttt{Jarmu} konstruktorát, beállítja az alap fejet (\texttt{SoproFej}), majd nullázza a készleteket (\texttt{so := 0}, \texttt{kerozin := 0}, \texttt{zuzottko := 0}).}
    \classitem{+statusLine(state: GameState): String \textit{(Override)}}{Részletes állapotleírást ad: pozíció, aktív fej, készletek és mozgásállapot.}
    \classitem{+name(): String \textit{(örökölt)}}{Visszaadja az azonosító nevet.}
    \classitem{+renameTo(newName: String): void \textit{(örökölt)}}{Frissíti az objektum nevét.}
    \classitem{+canMove(): boolean \textit{(örökölt)}}{Igazat ad vissza, ha a hókotró mozgásképes (\texttt{disabledTime <= 0}).}
    \classitem{+tickTime(): void \textit{(örökölt)}}{Időfrissítéskor csökkenti a mozgásképtelen időt, ha szükséges.}
    \classitem{+aktualisUt(state: GameState): Ut}{Visszaadja a hókotró aktuális útját. Hibát dob, ha a hókotró nincs úton vagy az út nem létezik.}
    \classitem{+takaritSav(ut: Ut, savIndex: int, state: GameState): void}{A megadott úton és sávon végrehajtja a takarítást az aktív fejjel. Ha nincs aktív fej, alapértelmezett \texttt{SoproFej}-et hoz létre.}
    \classitem{+fejCsere(jatekos: Jatekos, ujFej: FejTipus): void}{Lecseréli az aktív fejet a játékos raktárából. Hibát dob, ha az új fej nincs készleten; a korábbi aktív fejet visszateszi a raktárba.}
    \classitem{+sotoltes(jatekos: Jatekos): void}{A játékostól levonja a sótöltés árát (50), majd \texttt{so := 100}. Hibát dob, ha nincs elég pénz.}
    \classitem{+kerozintoltes(jatekos: Jatekos): void}{A játékostól levonja a kerozintöltés árát (60), majd \texttt{kerozin := 100}. Hibát dob, ha nincs elég pénz.}
    \classitem{+zuzalektoltes(jatekos: Jatekos): void}{A játékostól levonja a zúzaléktöltés árát (40), majd \texttt{zuzottko := 100}. Hibát dob, ha nincs elég pénz.}
\end{class-template-method}

\textbf{Pszeudokód a takarítás menetéhez:}
\begin{verbatim}
METHOD Takarit(hokotroNev, savIndex)
    IF hokotro.currentUt IS null THEN
     PRINT "A hokotro nincs uton"
    END IF

    ut := FindUt(hokotro.currentUt)
    sav := ut.GetSav(savIndex)

    IF hokotro.aktivFej IS null THEN
     hokotro.aktivFej := CreateFej(SOPROFEJ)
    END IF

    hokotro.aktivFej.ApplyCleaningEffect(hokotro, sav, ut, savIndex, state)
    LogEvent("A takaritas vegrehajtva")
END METHOD
\end{verbatim}

\textbf{Pszeudokód a hókotró fejének cseréjéhez:}
\begin{verbatim}
METHOD FejCsere(hokotroNev, ujFejTipus)
 IF jatekos NOT atTelephely THEN
  PRINT "Fejcsere csak telephelyen engedelyezett"
 END IF

 ujFej := FindFejInInventory(jatekos, ujFejTipus)
 IF ujFej IS null THEN
  PRINT "A kivant fej nincs a raktarban"
 END IF

 regiFej := hokotro.aktivFej
 hokotro.aktivFej := ujFej
 RemoveFejFromInventory(jatekos, ujFej)
 IF regiFej IS NOT null THEN
  AddFejToInventory(jatekos, regiFej.tipus())
 END IF

 LogEvent("A fejcsere sikeres")
END METHOD
\end{verbatim}


\subsection{Fej}
\begin{class-template-responsibility}
    Absztrakt ősosztály a hókotrókra szerelhető különböző takarítóeszközök számára. Definiálja az egységes takarítási interfészt. 
\end{class-template-responsibility}
\begin{class-template-attribute}
    \classitem{-tipus: FejTipus}{A fej típusa, amely meghatározza a takarítási viselkedést.}
\end{class-template-attribute}
\begin{class-template-method}
    \classitem{#Fej(tipus: FejTipus)}{Protected konstruktor. Beállítja a fej típusát.}
    \classitem{+tipus(): FejTipus}{Visszaadja az aktuális fej típusát.}
    \classitem{+takaritHatas(h: Hokotro, sav: Sav, ut: Ut, savIndex: int, state: GameState): void \textit{(abstract)}}{Absztrakt metódus. A konkrét fejek ebben adják meg a saját takarítási logikájukat.}
\end{class-template-method}


\subsection{SoproFej}
\begin{class-template-responsibility}
    Az alapértelmezett takarítófej, amely a csapadékot jobbra eső szomszédos sávba tolja át.
\end{class-template-responsibility}
\begin{class-template-baseclass}
    Fej \baseclass SoproFej
\end{class-template-baseclass}
\begin{class-template-attribute}
    \classitem{(nincs saját attribútum)}{Nincs külön állapota.}
\end{class-template-attribute}
\begin{class-template-method}
    \classitem{-SoproFej()}{Privát konstruktor. A fej típusát \texttt{SOPROFEJ}-re állítja.}
    \classitem{+takaritHatas(h: Hokotro, sav: Sav, ut: Ut, savIndex: int, state: GameState): void \textit{(Override)}}{A sávon lévő hó és feltört jég mennyiségét összegyűjti, a sávot lenullázza, majd ha van jobb oldali sáv, oda átteszi az összegyűjtött mennyiséget. Eseményt naplóz.}
\end{class-template-method}
\textbf{Pszeudokód a SoproFej működéséhez:}
\begin{verbatim}
METHOD: SoproFej.takaritHatas(h, sav, ut, savIndex, state)
    movedHo := sav.ho + sav.feltortJeg
    sav.ho := 0
    sav.feltortJeg := 0
    
    rightIndex := savIndex + 1
    IF movedHo > 0 AND rightIndex < ut.savSzam() THEN
        rightSav := ut.sav(rightIndex)
        rightSav.ho := rightSav.ho + movedHo
    END IF

    LogEvent("SoproFej takaritas vegrehajtva: " + movedSnow + " ho es feltort jeg attolve a jobb oldali savba, ha letezik.")
END METHOD
\end{verbatim}

\subsection{HanyoFej}
\begin{class-template-responsibility}
    Nagyteljesítményű fej, amely a hókotró nyomvonalából a havat és a feltört jeget a pályáról véglegesen eltávolítja. 
\end{class-template-responsibility}
\begin{class-template-baseclass}
    Fej \baseclass HanyoFej
\end{class-template-baseclass}
\begin{class-template-attribute}
    \classitem{(nincs saját attribútum)}{Nincs külön állapota.}
\end{class-template-attribute}
\begin{class-template-method}
    \classitem{-HanyoFej()}{Privát konstruktor. A fej típusát \texttt{HANYOFEJ}-re állítja.}
    \classitem{+takaritHatas(h: Hokotro, sav: Sav, ut: Ut, savIndex: int, state: GameState): void \textit{(Override)}}{A sávon lévő hó és feltört jég mennyiségét nullázza, majd takarítási eseményt vesz fel.}
\end{class-template-method}

\textbf{Pszeudokód a HanyoFej működéséhez:}
\begin{verbatim}
METHOD: HanyoFej.takaritHatas(h, sav, ut, savIndex, state)
    sav.ho := 0
    sav.feltortJeg := 0
    LogEvent("Hanyofej az ut melle tavolitotta a csapadekot.")
END METHOD
\end{verbatim}

\subsection{JegtoroFej}
\begin{class-template-responsibility}
    A masszív jégpáncél fizikai feltörését végző fej. 
\end{class-template-responsibility}
\begin{class-template-baseclass}
    Fej \baseclass JegtoroFej
\end{class-template-baseclass}
\begin{class-template-attribute}
    \classitem{(nincs saját attribútum)}{Nincs külön állapota.}
\end{class-template-attribute}
\begin{class-template-method}
    \classitem{-JegtoroFej()}{Privát konstruktor. A fej típusát \texttt{JEGTOROFEJ}-re állítja.}
    \classitem{+takaritHatas(h: Hokotro, sav: Sav, ut: Ut, savIndex: int, state: GameState): void \textit{(Override)}}{A jégréteget feltöri, a jégből feltört jégréteget hoz létre, majd ezt hóként kezeli a sávon. Eseményt naplóz.}
\end{class-template-method}

\textbf{Pszeudokód a JegtoroFej működéséhez:}
\begin{verbatim}
METHOD: JegtoroFej.takaritHatas(h, sav, ut, savIndex, state)
    sav.feltortJeg := sav.feltortJeg + sav.jeg
    sav.jeg := 0
    sav.ho := sav.ho + sav.feltortJeg
    sav.feltortJeg := 0
    LogEvent("Jegtorofej feltorte a jeget, ho jellegu retegre valtott.")
END METHOD
\end{verbatim}

\subsection{SoszoroFej}
\begin{class-template-responsibility}
    Só szórásával kémiai úton olvasztja el az akadályokat. Egyszerre egy sávot tud sózni. 
\end{class-template-responsibility}
\begin{class-template-baseclass}
    Fejj \baseclass SoszoroFej
\end{class-template-baseclass}
\begin{class-template-attribute}
    \classitem{(nincs saját attribútum)}{Nincs külön állapota.}
\end{class-template-attribute}
\begin{class-template-method}
    \classitem{-SoszoroFej()}{Privát konstruktor. A fej típusát \texttt{SOSZOROFEJ}-re állítja.}
    \classitem{+takaritHatas(h: Hokotro, sav: Sav, ut: Ut, savIndex: int, state: GameState): void \textit{(Override)}}{10 egység sót fogyaszt, a sáv hó- és jégkészletét lenullázza, majd 3 körre sózott védelmet állít be.}
\end{class-template-method}

\textbf{Pszeudokód a SoszoroFej működéséhez:}
\begin{verbatim}
METHOD: SoszoroFej.takaritHatas(h, sav, ut, savIndex, state)
    IF h.so < 10 THEN
        PRINT "Nincs eleg so a soszorofejhez."
        RETURN
    END IF
    h.so := h.so -10
    sav.jeg := 0
    sav.ho := 0
    sav.saltedRounds := 3
    LogEvent("Soszoro fej aktiv: ho es jeg eltunt, 3 korig vedett a sav.")
END METHOD
\end{verbatim}

\subsection{SarkanyFej}
\begin{class-template-responsibility}
    Prémium takarítófej, amely biokerozint égetve azonnali jég- és hóolvasztást végez és kerozint fogyaszt.  
\end{class-template-responsibility}
\begin{class-template-baseclass}
    Fej \baseclass SarkanyFej
\end{class-template-baseclass}
\begin{class-template-attribute}
    \classitem{(nincs saját attribútum)}{Nincs külön állapota.}
\end{class-template-attribute}
\begin{class-template-method}
    \classitem{-SarkanyFej()}{Privát konstruktor. A fej típusát \texttt{SARKANYFEJ}-re állítja.}
    \classitem{+takaritHatas(h: Hokotro, sav: Sav, ut: Ut, savIndex: int, state: GameState): void \textit{(Override)}}{10 egység kerozint fogyaszt, és a sáv minden csapadékalapú lerakódását eltávolítja.}
\end{class-template-method}

\textbf{Pszeudokód a SarkanyFej működéséhez:}
\begin{verbatim}
METHOD: SoszoroFej.takaritHatas(h, sav, ut, savIndex, state)
    IF h.kerozin < 10 THEN
        PRINT "Nincs eleg kerozin a sarkanyfejhez."
        RETURN
    END IF
    h.kerozin := h.so -10
    sav.jeg := 0
    sav.ho := 0
    sav.feltortJeg := 0
    LogEvent("Sarkanyfej aktiv: azonnali olvasztas vegrehajtva.")
END METHOD
\end{verbatim}

\subsection{ZuzottkoSzoroFej}
\begin{class-template-responsibility}
A    rendkívül veszélyes jégpáncél csúszósságának a megszűntetésére szolgáló fej, amely a sávra szórja a fejben tárolt zúzottkövet. Fontos felelőssége, hogy a sávon lévő csapadékot nem távolítja el.
\end{class-template-responsibility}
\begin{class-template-baseclass}
    Fej \baseclass ZuzottkoSzoroFej
\end{class-template-baseclass}
\begin{class-template-attribute}
    \classitem{(nincs saját attribútum)}{Nincs külön állapota.}
\end{class-template-attribute}
\begin{class-template-method}
    \classitem{-ZuzottkoSzoroFej()}{Privát konstruktor. A fej típusát \texttt{ZUZOTTKOSZOROFEJ}-re állítja.}
    \classitem{+takaritHatas(h: Hokotro, sav: Sav, ut: Ut, savIndex: int, state: GameState): void \textit{(Override)}}{10 egység zúzottkövet fogyaszt, majd 3 körre csúszásvédelmet állít be a sávra.}
\end{class-template-method}

\textbf{Pszeudokód a SarkanyFej működéséhez:}
\begin{verbatim}
METHOD: SoszoroFej.takaritHatas(h, sav, ut, savIndex, state)
    IF h.kerozin < 10 THEN
        PRINT "Nincs eleg zuzottko a zuzottkoszorofejhez."
        RETURN
    END IF
    h.zuzottko := h.zuzottko -10
    sav.grittedRounds := 3
    LogEvent("Zuzottko kiszorva: csuszasveszely csokkent 3 korre.")
END METHOD
\end{verbatim}

\subsection{Ut}
\begin{class-template-responsibility}
    Az úthálózat egy élét (szakaszát) reprezentáló osztály, amely összeköt két csomópontot. Tartalmazza az úton kialakított sávokat, tárolja az út típusát (normál, híd, alagút) és hosszát. Felelős az irányítással kapcsolatos alapvető gráf-információk szolgáltatásáért (pl. egyik végpont alapján a másik megállapítása), valamint a rajta lévő sávok összefogásáért.
\end{class-template-responsibility}
\begin{class-template-baseclass}
    Object \baseclass Ut
\end{class-template-baseclass}
\begin{class-template-attribute}
    \classitem{-name: String}{Az út egyedi azonosítója (neve).}
    \classitem{-nodeA: String}{Az utat határoló egyik csomópont neve.}
    \classitem{-nodeB: String}{Az utat határoló másik csomópont neve.}
    \classitem{-length: int}{Az út hossza.}
    \classitem{-type: UtTipus}{Az út típusa (\texttt{NORMAL}, \texttt{HID}, \texttt{ALAGUT}).}
    \classitem{-savok: List<Sav>}{Az úthoz tartozó sávok listája, melyek az akadályokat, lerakódásokat tárolják.}
\end{class-template-attribute}
\begin{class-template-method}
    \classitem{-Ut(name: String, nodeA: String, nodeB: String, length: int, type: UtTipus, savSzam: int)}{Privát konstruktor, inicializálja az utat és létrehozza a megadott számú sávot.}
    \classitem{+name(): String}{Visszaadja az út azonosító nevét.}
    \classitem{+renameTo(newName: String): void}{Segítségével megváltoztatható az út azonosítója a megfelelő esemény hatására.}
    \classitem{+hasNode(node: String): boolean}{Ellenőrzi és visszaadja, hogy a paraméterként kapott csomópont az út valamelyik végpontja-e.}
    \classitem{+opposite(node: String): String}{Visszaadja az út másik végpontját a paraméterként kapott (ismert) végponthoz képest, vagy \texttt{null}-t, ha az út nem tartalmazza a megadott végpontot.}
    \classitem{+savSzam(): int}{Visszaadja az úthoz tartozó sávok darabszámát.}
    \classitem{+sav(idx: int): Sav}{Visszaadja az úthoz tartozó felparaméterezett, adott sávindexeddel rendelkező sávot.}
    \classitem{+sumHo(): int}{Összeadja az út összes sávjában található hó mennyiségét, és visszaadja az eredményt.}
    \classitem{+alkalmazHofall(): int}{Minden sávon alkalmaz egy havazási lépést (alagút-védelem figyelembevételével), majd visszaadja a hó mennyiségének nettó növekményét (\texttt{after - before}).}
    \classitem{+describeCompact(): String}{Tömör, olvasható állapotleírást biztosít a konzolos kimenetre az útról, szomszédairól és sávjairól.}
\end{class-template-method}

\subsection{Sav}
\begin{class-template-responsibility}
    Egy út (Ut) adott sávját reprezentáló osztály. Felelős az adott sávon történő időjárás okozta lerakódások (hó, jég, feltört jég) tárolásáért, valamint a takarítási védelem (sózás, zúzottkő) hátralévő idejének nyilvántartásáért. Ő kezeli a hóesés hatásait és a védelem körönkénti csökkenését (idő telése).
\end{class-template-responsibility}
\begin{class-template-baseclass}
    Object \baseclass Sav
\end{class-template-baseclass}
\begin{class-template-attribute}
    \classitem{-index: int}{A sáv indexe (azonosítója) a befogadó úton belül.}
    \classitem{-snow: int}{A sávon felgyülemlett hó aktuális mennyisége.}
    \classitem{-ice: int}{A sávon lévő jég aktuális mennyisége.}
    \classitem{-brokenIce: int}{A sávon lévő, jégtörő fejjel már feltört jég mennyisége.}
    \classitem{-saltedRounds: int}{A hátralévő körök száma, ameddig a sáv sózva van (ezalatt a ráhulló hó/jég elolvad).}
    \classitem{-grittedRounds: int}{A hátralévő körök száma, ameddig a sáv zúzottkővel szórt (csökkenti a csúszásveszélyt).}
\end{class-template-attribute}
\begin{class-template-method}
    \classitem{-Sav(index: int)}{Privát konstruktor. Beállítja a sáv indexét, a lerakódásokat és a védelmeket alapesetben nullára inicializálva.}
    \classitem{+applyDeposit(type: DepositType, value: int): void}{Célzottan helyez el a sávon különböző típusú lerakódásokat (\texttt{HO}, \texttt{JEG}, \texttt{SO}, \texttt{ZUZALEK}, \texttt{FELTORTJEG}) a megadott mennyiségben vagy értéken.}
    \classitem{+snowingOneUnit(protectedByTunnel: boolean): void}{Egy egységnyi havat próbál a sávhoz adni havazás esemény esetén, figyelembe véve, hogy az út alagútban van-e (védve a havazástól), vagy éppen sózva van-e a sáv.}
    \classitem{+tickTime(): void}{Körléptetés. Eggyel csökkenti a sózott és a zúzottköves állapotból hátralévő körök számát, ha azok nagyobbak nullánál.}
    \classitem{+describeCompact(): String}{Tömör, kifejező karakteres leírást nyújt a sáv indexéről, hószintjéről, jégszintjéről és egyedi flaggel a sózottságáról a konzolos kimenet számára.}
\end{class-template-method}

\subsection{Jatekos}
\begin{class-template-responsibility}
    A menedzser-játékosokat reprezentáló osztály, amely kezeli a pénzügyeket, járműlistát és fejraktárt. Refaktor után ide került a vásárlási logika egy része.
\end{class-template-responsibility}
\begin{class-template-baseclass}
    Object \baseclass Jatekos
\end{class-template-baseclass}
\begin{class-template-attribute}
    \classitem{\#name: String}{A játékos neve/azonosítója.}
    \classitem{\#money: int}{A játékos aktuális pénze.}
    \classitem{\#vehicles: List<String>}{A játékoshoz tartozó járművek neveinek listája.}
    \classitem{\#inventory: List<FejTipus>}{A játékos raktárában elérhető fejtípusok listája.}
\end{class-template-attribute}
\begin{class-template-method}
    \classitem{+vasarolHokotro(hokotroNev: String): Hokotro}{Levásárolja a hókotró árát (300), létrehozza az új hókotrót, hozzárendeli a tulajdonost és a játékos járműlistájához adja. Hibát dob, ha nincs elég pénz.}
    \classitem{+vasarolFej(fejTipus: FejTipus): void}{Levásárolja a fej árát (80), majd a választott fejtípust hozzáadja a játékos raktárához. Hibát dob, ha nincs elég pénz.}
\end{class-template-method}

\section{A tesztek részletes tervei, leírásuk a teszt nyelvén}





\section{A tesztek részletes tervei, leírásuk a teszt nyelvén}
\comment{A tesztek részletes tervei alatt meg kell adni azokat a bemeneti adatsorozatokat, amelyekkel a program működése ellenőrizhető. Minden bemenő adatsorozathoz definiálni kell, hogy az adatsorozat végrehajtásától a program mely részeinek, funkcióinak ellenőrzését várjuk és konkrétan milyen eredményekre számítunk, ezek az eredmények hogyan vethetők össze a bemenetekkel.A tesztek leírásakor az előző dokumentumban (proto koncepciója) megadott szintakszist kell használni.}

\subsection{Teszteset1}
\begin{test-case-description}
    Szöveges leírás, kb 1-5 mondat
\end{test-case-description}
\begin{test-case-function}
    ...
\end{test-case-function}
\begin{test-case-input}
    A proto nyelvén megadva, pl.:
    \begin{verbatim}
    init world
    0: S 3
    ...
    print
    \end{verbatim}
\end{test-case-input}
\begin{test-case-output}
    A poro kimeneti nyelvén megadva, pl.:
    \begin{verbatim}
    worlddata
    ...
    creatures 0
    \end{verbatim}
\end{test-case-output}

\section{A tesztelést támogató programok tervei}
\comment{A tesztadatok előállítására, a tesztek eredményeinek kiértékelésére szolgáló segédprogramok részletes terveit kell elkészíteni.}